% Section 2 — PSM architecture
% Status: v1 draft. Describes the engine's three data structures
% (H3 cells / ring buffer of HLLs / reservoir-sampled exemplars)
% and the per_cell_cap selection rule. Anchored on the actual
% C-side surface in include/core/{spatial_memory,tile,ring_buffer}.h
% so any future reader can navigate from the prose to the code.
% NB: former §3 (PSM-as-MLLM-prefilter pipeline) folded in as the final
% subsection, labelled sec:pipeline, during the 8pp cut.

\section{Probabilistic Spatial Memory}

\label{sec:method}

\psm{} is a streaming spatial-temporal index that ingests
$(t, \mathrm{lat}, \mathrm{lng}, \mathbf{e})$ tuples: timestamp,
geographic position, and a $d$-dimensional CLIP embedding. This supports
top-$K$ similarity search with capped per-cell state at any time. Three
components define the index.



\subsection{H3 spatial tessellation}
\label{sec:method-h3}

We tessellate the wearer's trajectory with H3~\cite{brodsky2018h3}, a
hierarchical hexagonal geospatial index. Each tuple maps to its containing cell
at resolution~$r$ (default~12, ${\sim}9$\,m edges, room scale; swept over
$r\in\{8,9,10,11,12\}$ in \S\ref{sec:results-multi-corpus}). H3 cell IDs are
global, comparable across recordings, sessions, and devices.
Positions come from whatever pose stream the device already runs: GNSS outdoors
(HDD's RTK fixes) or geo-anchored on-device SLAM/VIO indoors (Aria MPS,
LiDAR-SLAM); \psm{} needs only a consistent global frame, so the indoor ``keys''
setting uses the same localization the device already provides.

\subsection{Time-decayed ring buffer of HyperLogLog sketches}
\label{sec:method-ring-buffer}

Each cell maintains a ring buffer of $C$ HyperLogLog
sketches~\cite{flajolet2007hyperloglog} rotating on a time-window cadence
(default 30\,s): ingest hashes the embedding into the current sketch, and when
the window elapses the oldest sketch expires as the ring wraps. The sketches
give per-cell cardinality and expose the cell's \emph{temporal envelope}
$(t_{\min},t_{\max})$, returned as the bucket window the MLLM consumes as a
time-range hint. Each sketch is $2^p$ bytes ($p{=}10\Rightarrow$1\,KiB), so the
ring costs $C\cdot1\,\mathrm{KiB}=60\,\mathrm{KiB}$ at $C{=}60$ per cell,
independent of dwell time. What the sketch counts is the number of \emph{distinct
hashed embedding vectors}: for continuous CLIP embeddings, where near-duplicate
frames still differ in their raw bytes, this is close to the ingested frame count
rather than a count of semantically distinct events or objects. We therefore use
the cardinality as a cheap per-cell activity/temporal signal and \textbf{do not}
claim it counts ``unique events'' or ``distinct people'' without an upstream
object/event canonicalization stage (\S\ref{sec:limitations}). The HLL is thus
optional metadata unless it is wired to drive allocation, which we leave to future
work.

\subsection{Reservoir-sampled exemplars}
\label{sec:method-exemplars}

Counts and envelopes are not embeddings, so each cell also stores exemplars via
\emph{Algorithm~R reservoir sampling}~\cite{vitter1985reservoir} (capacity
$R{=}128$): after $n>R$ frames the reservoir holds a uniform random $R$-subset
(${\sim}0.38$\,MiB/cell at $R{=}128$, CLIP-L's $d{=}768$; ${\sim}0.44$\,MiB with
the HLL ring).

\paragraph{Memory profile (per-cell bounded, not globally bounded).} For a
trajectory visiting $C_n$ distinct H3 cells over $n$ frames, \psm{}'s state is
$\Theta\!\big(C_n(R d + C\,2^p)\big)$, whereas brute-force embedding storage is
$\Theta(n d)$. The \emph{per-cell} summary is bounded ($R$ exemplars plus a
$C$-slot HLL ring, independent of how long the wearer dwells), so state grows
sublinearly in $n$ \emph{when revisits dominate}: revisited and dwelled cells
saturate at $R$ rather than accumulating every frame. This is not a worst-case
bound independent of recording duration: $C_n$ grows with the physical area
explored, and under continual exploration $C_n=\Theta(n)$, so total state is
unbounded in the worst case. The HLL ring is charged to \emph{every} visited
cell, including the many single-visit ones, so on wide-area corpora it can
dominate the footprint (\S\ref{sec:results-hdd}). We therefore report exemplar
and ring bytes separately and, in \S\ref{sec:results}, study a \emph{fixed
global exemplar budget} that allocates a hard exemplar-payload ceiling across
cells; per-cell policy bookkeeping is reported separately.
Figure~\ref{fig:f8-psm-viz} shows
the per-cell structures over a real trajectory.

\paragraph{Exemplar/window lifetimes.} Exemplars persist for the cell's lifetime
(reservoir sampling only \emph{replaces} slots; it has no time-based expiry),
while the HLL ring rotates its oldest window out on the time cadence. A retained
exemplar's timestamp can therefore outlive every HLL window covering it; the
returned bucket envelope $(t_{\min},t_{\max})$ reflects the current ring, not the
exemplar's original interval. Aligning exemplar aging with the window (decay-aware
sampling) is future work.

\subsection{Top-$K$ query with \caponek}
\label{sec:method-query}

A query takes a vector $\mathbf{q}$ and, optionally, a geographic center and
k-ring radius. Each in-scope cell returns at most \caponek{} exemplars ranked by
cosine similarity; the engine sorts their union and emits the top $K$ with each
cell's bucket window. $\caponek{=}1$ enforces strict spatial diversity
(one slot per cell), $\caponek{=}K$ recovers plain top-$K$, and intermediate
values trade off smoothly (\S\ref{sec:results-engine}). Ingest is
$O(1)$ amortized. Query cost depends on scope: a \emph{centered} query supplies an
$(\mathrm{lat},\mathrm{lng})$ center with a k-ring radius, so only the in-scope
cells are scanned; an \emph{uncentered} (global) query has no spatial prefilter
and scores every exemplar in every cell, i.e.\ $\Theta$(total retained
exemplars) per query. Our target-configuration host benchmark measures the
uncentered global-scan path (\S\ref{sec:results-memory-latency}); scaling it via
per-cell prototypes or another router is future work.

\subsection{PSM as an MLLM prefilter}
\label{sec:pipeline}
\psm{} returns ranked tuples, not answers, so we compose it as a
\emph{prefilter} feeding a frontier MLLM \emph{reranker}
(Fig.~\ref{fig:f1-architecture}, bottom row): embed the question, call
\texttt{query\_similar} at the configured \caponek{} and $K{=}5$, decode the five
candidate frames on demand (a device-cache memcpy in deployment, not a video
seek), and pass them with $q$ in one Gemini~3.1~Pro~\cite{google2026gemini} call
whose frozen prompt selects the best frame. The split is deliberate: \psm{}
retrieves candidates and the MLLM selects one (\S\ref{sec:results-engine}).
Encoder choice is a pipeline knob, not an
architectural commitment: one factory serves CLIP-L, CLIP-bigG, and
SigLIP-2-large~\cite{tschannen2025siglip2}\footnote{Encoder scales: CLIP-L
430\,M/768\,d, CLIP-bigG 2.5\,B/1280\,d, SigLIP-2-large 660\,M/1024\,d.}; the C
engine sees only a $D$-dim vector. We evaluate resolution sensitivity with all
three encoders (\S\ref{sec:results-multi-corpus}).
